%---------------------------------------------------------------------------%
%->> 封面信息及生成
%---------------------------------------------------------------------------%
%-
%-> 中文封面信息（按用户要求，个人信息留空）
%-
\confidential{}% 密级：只有涉密论文才填写
\schoollogo{scale=0.2}{ShanghaiTechLogo}% 校徽
\title{基于异构图与多智能体协商的家谱婚姻链接重建可视分析}% 论文中文题目
\author{}% 论文作者：留空
\ID{}
\entranceYear{}
\advisor{}% 指导教师：留空
\advisorsec{}% 第二指导老师
\degree{学士}% 学位
\degreetype{工学}% 学位类别
\major{计算机科学与技术}% 二级学科专业名称
\institute{上海科技大学信息科学与技术学院}% 院系名称
\chinesedate{二零二六年六月}% 毕业日期占位
%-
%-> 英文封面信息
%-
\englishtitle{Visual Analytics for Reconstructing Marriage Edges in Patrilineal Genealogies via Heterogeneous Graph Transformer and Multi-Agent Negotiation}
\englishauthor{}
\englishadvisor{}
\englishdegree{Bachelor}
\englishdegreetype{Engineering}
\englishthesistype{thesis}
\englishmajor{Computer Science and Technology}
\englishinstitute{School of Information Science and Technology, ShanghaiTech University}
\englishdate{June 2026}
%-
%-======================================================
% 封面和声明页格式不符合教务处要求，因此此处不生成pdf格式的 ||
% 论文封面和声明页，请各位使用word模板中的相关内容。       ||
%=======================================================
\maketitle% 生成中文封面
\makeenglishtitle% 生成英文封面
%-
%-> 作者声明
%-
\makedeclaration% 生成声明页
%-
%-> 中文摘要
%-
\setcounter{page}{1}
\pagenumbering{Roman}
\chapter*{摘\quad 要}\chaptermark{摘\quad 要}% 摘要标题

家谱是研究中国前现代人口史与社会史的核心一手史料，但因父系记录传统的限制，跨宗族的婚姻关系与母系世系长期被系统性地遮蔽，使原本完整的亲属网络在文献层面碎裂为彼此孤立的家族树。如何在缺失婚姻边的家谱图谱中识别可信的婚姻候选，并使预测过程对历史学者可解释、可干预、可回溯，是计算社会史与可视分析交叉研究中尚未很好解决的问题。本文以中国多代人口面板数据集（CMGPD-LN）为实验载体，提出一套将异构图变换器（HGT）链路预测、SEAL 子图模式提取与基于大语言模型的多智能体协商相耦合的家谱婚姻边重建可视分析框架\textsc{GeneaLink}。首先，对原始记录构建包含 4 类节点、9 类边的异构图，并通过对母子边的"消融而不删除"实现对真实家谱缺失模式的镜像；随后训练一个二层、四头注意力、隐藏维度 128 的 HGT 编码器，搭配以拼接、绝对差与哈达玛积为输入的婚姻评分多层感知机；之后利用 SEAL 框架提取候选对的封闭子图，并应用双半径节点标注得到 13 类与历史学语义相关的微结构图样；最后将上述结构证据与档案事件字段共同注入两个候选人侧的智能体角色中，让其在六轮预设流程中以"印象生成—深入询问—质疑—对齐—复述—终评"的方式协商，得到对称惩罚下的最终评分。在此之上，论文进一步设计了一套包含六个联动视图的混合主动可视分析界面：性能概览、流程视图、嵌入视图（蜂窝画布）、二分图视图、协商竞技场视图与规则视图，使专家既能批量接受高置信对，又能对边界案例逐对审视、注入领域知识、回滚错误决策。专家研究表明，该系统在可用性与可解释性上均获得 6.0/7 的均值评分，提示注入与来源标记成为最被重视的两项功能。本工作的方法论贡献在于把图学习、子图模式与智能体推理之间的鸿沟通过可视化进行有机连接，为复杂社会网络的链路预测可视分析提供了一条可复用的研究路径。

\keywords{异构图变换器，子图链路预测，多智能体模拟，家谱重建，可视分析}
%-
%-> 英文摘要
%-
\chapter*{Abstract}\chaptermark{Abstract}

Chinese patrilineal genealogies are a central primary source for pre-modern demographic and social history, yet by recording convention they systematically conceal cross-clan marital ties and maternal lineage, fragmenting the underlying kinship network into isolated family trees. Restoring the missing marriage edges in such fragmented graphs in a way that historians can inspect, steer, and revoke remains an open problem at the intersection of computational social history and visual analytics. This thesis presents \textsc{GeneaLink}, a visual-analytics framework that couples Heterogeneous Graph Transformer (HGT) link prediction, SEAL subgraph motif extraction, and large-language-model multi-agent negotiation, applied to the China Multi-Generational Panel Dataset for Liaoning (CMGPD-LN). The pipeline (i) builds a heterogeneous graph with four node types and nine edge types and ablates maternal edges in place to mirror real-world genealogical deficit; (ii) trains a two-layer, four-head HGT encoder of hidden width 128 together with a marriage scorer that concatenates pair embeddings, their absolute difference, and their Hadamard product; (iii) extracts thirteen domain-relevant motifs from candidate-enclosing subgraphs via Double-Radius Node Labeling; and (iv) instantiates two persona agents grounded in archival records that negotiate over six structured rounds with a symmetry-penalised final score. A six-view interface---performance overview, process ledger, honeycomb embedding canvas, bipartite acceptance board, agent arena, and rule-transfer panel---supports batch acceptance, single-pair adjudication, hint injection, and one-click rollback. Expert evaluation with five domain specialists yielded mean Likert ratings of 6.0/7 for usefulness and interpretability and identified the hint console and provenance chips as the most appreciated features. The contribution of this work lies in bridging graph learning, subgraph motifs, and agent reasoning through tightly coupled visualisation, offering a reusable methodology for link-prediction visual analytics over complex social networks.

\englishkeywords{Heterogeneous Graph Transformer, SEAL link prediction, Multi-Agent Simulation, genealogy reconstruction, Visual Analytics}
%---------------------------------------------------------------------------%
